\documentclass[12pt,a4paper]{article}

% --- Packages de base ---
\usepackage[utf8]{inputenc}
\usepackage[T1]{fontenc}
\usepackage[french]{babel}
\usepackage{geometry}
\geometry{margin=2cm, left=3cm} 
\usepackage{xcolor}
\usepackage{titlesec}
\usepackage{array}
\usepackage{colortbl}
\usepackage{enumitem}
\usepackage{fontawesome5}
\usepackage{graphicx}
\usepackage{tcolorbox}

% --- Couleurs Officielles ---
\definecolor{GrillRed}{HTML}{D32F2F}
\definecolor{CoalBlack}{HTML}{2D2D2D}
\definecolor{SauceCream}{HTML}{FDFBF7}
\definecolor{FryGold}{HTML}{FFC107}
\definecolor{StoneGray}{HTML}{BDBDBD}
\definecolor{DarkRed}{HTML}{B71C1C} % Pour le hover
\definecolor{BrownVisited}{HTML}{5D4037} % Pour les liens visités
% Nouvelles couleurs sémantiques pour les boutons et statuts
\definecolor{SuccessGreen}{HTML}{2E7D32}
\definecolor{GpsBlue}{HTML}{1565C0}
\definecolor{HoverGray}{HTML}{9E9E9E}

% --- Configuration de la page ---
\pagecolor{white}
\setlength{\parindent}{0pt}

% Style des titres
\titleformat{\section}{\color{GrillRed}\normalfont\Large\bfseries\uppercase}{}{0em}{}[\titlerule]
\titleformat{\subsection}{\color{CoalBlack}\normalfont\large\bfseries}{}{0em}{\faCaretRight\quad}

% Bandeau latéral décoratif
\usepackage{background}
\backgroundsetup{
scale=1, angle=0, opacity=1, contents={
    \begin{tikzpicture}[remember picture, overlay]
        \path [fill=CoalBlack] (current page.north west) rectangle ([xshift=1cm]current page.south west);
        \path [fill=GrillRed] (current page.north west) rectangle ([xshift=0.2cm]current page.south west);
    \end{tikzpicture}}
}

\begin{document}

% --- HEADER ---
\begin{center}
    \includegraphics[width=5cm]{logo-le-grand-miam-removebg-preview.png} \\
    \vspace{0.4cm}
    {\Huge \textbf{\color{CoalBlack}DOCUMENT DE \color{GrillRed}CONCEPTION}} \\
    \vspace{0.2cm}
    \texttt{\small PROJET YUMLAND | LE GRAND MIAM} \\
    \vspace{0.2cm}
    \textit{\small Thème : Steakhouse \& Grillades Premium} \\
    \vspace{0.4cm}
    \textbf{Groupe MI2A} \\
    {\normalsize Myriam BENSAÏD \quad Sheryne OUARGHI-MHIRI \quad Kylian VANDEL} \\
    
\end{center}
\vspace{1cm}

\section{1. Identité Visuelle et Concept}

\textbf{Le Grand Miam} est un restaurant convivial spécialisé dans les viandes grillées et les portions généreuses. L'identité visuelle repose sur l'univers du "Grill" : un rendu à la fois chaleureux, robuste et appétissant.

\begin{tcolorbox}[colback=SauceCream, colframe=CoalBlack, arc=2mm, boxrule=1pt, drop shadow]
    \textbf{\large La Stratégie : "Le Terroir pour Tous"} \\
    \textbf{La Gestion des Viandes :}
    \begin{itemize}[label=\textcolor{GrillRed}{\faCheck}, leftmargin=*]
        \item \textbf{Bœuf :} Mise en avant de l'origine France (Charolais/Limousin). Il est indiqué clairement : \textit{"Toutes nos pièces de bœuf sont disponibles en version Halal sur demande."}
        \item \textbf{Porc :} Présent dans une section isolée "Spécialités Charcutières" pour bien séparer les produits.
    \end{itemize}
    \vspace{0.2cm}
    \textbf{L'interface est conçue pour être :}
    \begin{itemize}[label=\textcolor{GrillRed}{\faCheck}, leftmargin=*]
        \item \textbf{Gourmande :} Mise en avant de visuels de haute qualité.
        \item \textbf{Accessible :} Contrastes forts pour une lisibilité parfaite (client et livreur).
        \item \textbf{Efficace :} Parcours utilisateur simplifié pour une commande rapide.
    \end{itemize}
\end{tcolorbox}

\section{2. Palette de Couleurs}
L'ambiance chromatique joue sur le contraste entre la chaleur du grill et la sobriété du charbon.

\begin{center}
\renewcommand{\arraystretch}{1.6}
\begin{tabular}{|m{4.5cm}|c|>{\centering\arraybackslash}m{1.5cm}|m{7cm}|}
    \hline
    \rowcolor{CoalBlack} \textcolor{white}{\textbf{Nom}} & \textcolor{white}{\textbf{HEX}} & \textcolor{white}{\textbf{Aperçu}} & \textcolor{white}{\textbf{Usage Principal}} \\
    \hline
    \textbf{\textcolor{GrillRed}{\faCircle} Rouge "Grill"} & \#D32F2F & \cellcolor{GrillRed} & Couleur Primaire. Boutons (CTA), prix, titres, icônes. \\
    \hline
    \textbf{\textcolor{CoalBlack}{\faCircle} Noir "Charbon"} & \#2D2D2D & \cellcolor{CoalBlack} & Couleur Secondaire. Header, Footer, navigation, textes. \\
    \hline
    \textbf{\textcolor{SauceCream}{\faCircle} Crème "Sauce"} & \#FDFBF7 & \cellcolor{SauceCream} & Arrière-plan (Body). Réduit la fatigue visuelle. \\
    \hline
    \textbf{\textcolor{FryGold}{\faCircle} Or "Frites"} & \#FFC107 & \cellcolor{FryGold} & Accent. Étoiles, promotions, focus formulaires. \\
    \hline
    \textbf{\textcolor{StoneGray}{\faCircle} Gris "Pierre"} & \#BDBDBD & \cellcolor{StoneGray} & Éléments passifs. Bordures, placeholders, textes secondaires. \\
    \hline
    \multicolumn{4}{|c|}{\cellcolor{CoalBlack} \textcolor{white}{\textbf{Couleurs Sémantiques (Statuts \& Métiers)}}} \\
    \hline
    \textbf{\textcolor{SuccessGreen}{\faCircle} Vert "Succès"} & \#2E7D32 & \cellcolor{SuccessGreen} & Boutons "Valider/Livrée", alertes de succès. \\
    \hline
    \textbf{\textcolor{GpsBlue}{\faCircle} Bleu "GPS"} & \#1565C0 & \cellcolor{GpsBlue} & Bouton de navigation externe (Livreur). \\
    \hline
\end{tabular}
\end{center}

\section{3. Typographie (Google Fonts)}
Combinaison de polices libres de droits pour allier caractère et lisibilité.

\begin{tcolorbox}[colback=SauceCream, colframe=CoalBlack!20, arc=2mm, title=\color{CoalBlack}\textbf{Titres (H1, H2, H3) : OSWALD}, fonttitle=\small, drop fuzzy shadow]
    {\Large\bfseries Grand par le goût, géant par l'appétit.} \\
    \vspace{0.1cm}
    \small \textbf{Style :} Sans-serif, Condensé, Épais (Bold / 700). \\
    \textbf{Usage :} Titres de pages, noms des plats, slogans. Donne un aspect "enseigne" et "industriel" propre aux steakhouses.
\end{tcolorbox}

\begin{tcolorbox}[colback=white, colframe=CoalBlack!20, arc=2mm, title=\color{CoalBlack}\textbf{Corps de Texte : LATO} (Fallback : Roboto), fonttitle=\small, drop fuzzy shadow]
    {\large Toutes nos viandes sont sélectionnées avec soin auprès d'éleveurs locaux.} \\
    \vspace{0.1cm}
    \small \textbf{Style :} Sans-serif, Rond, Régulier (Regular / 400). \\
    \textbf{Usage :} Paragraphes, liens, inputs, mentions légales. \\
    \textbf{Taille de base :} 16px (Desktop) / 18px (Mobile \& Livreur).
\end{tcolorbox}

\section{4. Éléments d'Interface (UI Kit)}

\subsection{Boutons (Call to Action \& Sémantiques)}
L'interface utilise un système de boutons hiérarchisé pour guider l'utilisateur. Tous les boutons partagent des propriétés communes pour garantir une cohérence visuelle :
\begin{itemize}[label=\textcolor{CoalBlack}{\tiny\faCircle}, itemsep=0pt]
    \item \textbf{Typographie :} Majuscules (\texttt{text-transform: uppercase}), police \textit{Oswald} en gras (\texttt{font-weight: 700}).
    \item \textbf{Forme \& Volume :} Arrondi strict (\texttt{border-radius: 4px}) et ombre portée (\texttt{box-shadow}) pour donner du relief.
    \item \textbf{Interactions :} Au survol, effet de zoom (\texttt{transform: scale(1.02)}) et assombrissement du fond.
\end{itemize}

\vspace{0.3cm}
\textbf{1. Bouton Primaire (Action principale / CTA)}
\begin{center}
    \begin{tcolorbox}[hbox, colback=GrillRed, colframe=DarkRed, arc=4px, boxrule=1pt, left=15pt, right=15pt, top=5pt, bottom=5pt, drop fuzzy shadow]
        \textcolor{SauceCream}{\textbf{\faCheck \quad COMMANDER / ENREGISTRER}}
    \end{tcolorbox}
\end{center}
\textit{Usage :} Ajouter au panier, enregistrer le profil, valider un formulaire. \\
\textit{Style :} Fond Rouge "Grill" (\#D32F2F), texte Crème "Sauce" (\#FDFBF7). Survol : \#B71C1C.

\vspace{0.3cm}
\textbf{2. Bouton Secondaire (Annuler / Retour)}
\begin{center}
    \begin{tcolorbox}[hbox, colback=StoneGray, colframe=HoverGray, arc=4px, boxrule=1pt, left=15pt, right=15pt, top=5pt, bottom=5pt, drop fuzzy shadow]
        \textcolor{CoalBlack}{\textbf{\faTimes \quad ANNULER}}
    \end{tcolorbox}
\end{center}
\textit{Usage :} Annuler une édition (ex: profil client), fermer une fenêtre modale. \\
\textit{Style :} Fond Gris "Pierre" (\#BDBDBD), texte Noir "Charbon" (\#2D2D2D). Survol : \#9E9E9E.

\vspace{0.3cm}
\textbf{3. Boutons Spécifiques Métiers (Livreur)}
\begin{center}
    \begin{tcolorbox}[hbox, colback=SuccessGreen, colframe=SuccessGreen!80!black, arc=4px, boxrule=1pt, left=15pt, right=15pt, top=5pt, bottom=5pt, drop fuzzy shadow]
        \textcolor{white}{\textbf{\faCheckCircle \quad VALIDÉE}}
    \end{tcolorbox}
    \hspace{10pt}
    \begin{tcolorbox}[hbox, colback=GpsBlue, colframe=GpsBlue!80!black, arc=4px, boxrule=1pt, left=15pt, right=15pt, top=5pt, bottom=5pt, drop fuzzy shadow]
        \textcolor{white}{\textbf{\faMapMarkerAlt \quad GPS}}
    \end{tcolorbox}
\end{center}
\textit{Usage :} Validation des courses et géolocalisation express. Adaptés pour un usage extérieur. \\
\textit{Style :} Vert foncé (\#2E7D32) et Bleu foncé (\#1565C0). \textbf{Hauteur minimale imposée à 60px} pour la préhension avec des gants.

\subsection{Formulaires \& Liens}
\begin{itemize}[label=\textcolor{CoalBlack}{\tiny\faCircle}, itemsep=0pt]
    \item \textbf{Champs :} Fond blanc sur fond de page crème. Bordures \#BDBDBD au repos.
    \item \textbf{Focus Saisie :} Bordure \#FFC107 (Or) et lueur externe.
    \item \textbf{Liens :} \textbf{Défaut} (Souligné/gras \#2D2D2D) | \textbf{Visité} (Brun \#5D4037) | \textbf{Survol} (\#D32F2F sans soulignement).
\end{itemize}

\section{5. Ergonomie Spécifique : Profil Livreur \faMotorcycle}
Pour répondre aux contraintes du livreur (smartphone, gants, extérieur) :
\begin{tcolorbox}[colback=white, colframe=GrillRed, arc=0mm, boxrule=1.5pt, drop fuzzy shadow]
    \begin{itemize}[label=\textcolor{GrillRed}{\faHandPointer}]
        \item \textbf{Zone de Toucher (Touch Target) :} Hauteur minimale de \textbf{60px} pour tous les éléments interactifs.
        \item \textbf{Espacement :} Marges de \textbf{20px} minimum pour éviter les erreurs ("Fat finger").
        \item \textbf{Contraste :} Texte noir sur fond blanc ou bouton rouge vif pour une lisibilité maximale en plein soleil.
    \end{itemize}
\end{tcolorbox}

\section{6. Maquettes et Structure (Wireframes)}
\begin{itemize}[label=\textcolor{CoalBlack}{\faLayerGroup}]
    \item \textbf{Header :} Logo à gauche, Navigation centrée, Bouton "Connexion/Panier" à droite.
    \item \textbf{Cartes Produits :} Visuel du plat en haut, Nom et description au centre, Prix et bouton d'ajout en bas. Effet d'ombre portée pour simuler une assiette posée sur une table.
    \item \textbf{Navigation Mobile :} Menu "Burger" simplifié pour les petits écrans.
\end{itemize}

\section{7. Accessibilité et Inclusivité}

Afin de rendre l'application "Le Grand Miam" accessible à tous les publics et d'offrir un confort de navigation optimal, deux modes spécifiques complètent cette charte :

\subsection{Thème Sombre (Dark Mode)}
Pensé pour réduire la fatigue visuelle (notamment pour les livreurs en fin de service) et pour s'adapter aux préférences des utilisateurs.
\begin{itemize}[label=\textcolor{CoalBlack}{\tiny\faCircle}, itemsep=0pt]
    \item \textbf{Fonds et Surfaces :} Remplacement de la couleur \textit{Crème "Sauce"} par un noir absolu (\#000000) et des gris anthracites (\#121212, \#1A1A1A) pour les conteneurs et les cartes.
    \item \textbf{Typographie :} Utilisation d'un gris très clair (\#E0E0E0) au lieu du blanc pur pour le texte, afin de limiter l'éblouissement.
    \item \textbf{Identité Visuelle :} Le \textit{Rouge "Grill"} et l'\textit{Or "Frites"} restent identiques et servent de couleurs d'accentuation.
\end{itemize}

\subsection{Mode Dyslexique (Accessibilité Cognitive)}
Une interface alternative repensée pour faciliter la lecture des personnes atteintes de dyslexie.
\begin{tcolorbox}[colback=white, colframe=CoalBlack, arc=2mm, boxrule=1pt, drop fuzzy shadow]
    \begin{itemize}[label=\textcolor{GrillRed}{\faCheck}, leftmargin=*]
        \item \textbf{Polices adaptées :} Bascule temporaire sur une typographie ultra-lisible (comme \textit{OpenDyslexic}, ou \textit{Verdana}/\textit{Arial}), en remplacement des polices plus graphiques.
        \item \textbf{Espacement accru :} Augmentation systématique de l'espacement entre les lettres (0.5px) et entre les mots (2px).
        \item \textbf{Aération :} Interlignage élargi (à 1.6em minimum) pour éviter le chevauchement visuel des lignes de texte.
    \end{itemize}
\end{tcolorbox}

\section{8. Iconographie (FontAwesome)}
L'iconographie joue un rôle clé pour la navigation rapide, réduire la charge cognitive et aider au repérage visuel, particulièrement pour les profils métiers.
\begin{itemize}[label=\textcolor{GrillRed}{\faChevronRight}]
    \item \textbf{Style Général :} Utilisation de la bibliothèque \textbf{FontAwesome 6} (variante \textit{Solid} privilégiée).
    \item \textbf{Navigation \& Actions :} \faHome (Accueil), \faUser (Espace Client), \faShoppingCart (Panier/Commande), \faSignOutAlt (Déconnexion), \faPlusCircle (Ajouter).
    \item \textbf{Outils Métiers :} \faMotorcycle (Livreur - Courses), \faFireBurner (Cuisine - Préparation), \faShieldAlt (Administration).
    \item \textbf{Indicateurs de Statut :} \faCheckCircle (Succès/Validé), \faExclamationTriangle (Alerte/Erreur), \faClock (En attente).
\end{itemize}

\section{9. Grille, Espacements et Responsive Design}
La mise en page utilise un design fluide (Flexbox/Grid) qui s'adapte à tous les écrans tout en conservant une structure aérée.

\begin{tcolorbox}[colback=SauceCream, colframe=CoalBlack!20, arc=2mm, title=\color{CoalBlack}\textbf{Règles d'Espacement (Marges \& Gouttières)}, fonttitle=\small, drop fuzzy shadow]
    \begin{itemize}[label=\textcolor{CoalBlack}{\tiny\faCircle}, itemsep=0pt]
        \item \textbf{Gouttières (Gaps) :} Minimum 15px à 20px entre les éléments (cartes, boutons) pour éviter toute surcharge.
        \item \textbf{Paddings internes :} 20px pour les cartes et conteneurs pour garantir que le texte respire.
        \item \textbf{Arrondis (Border-radius) :} Style hiérarchisé : \textbf{12px} pour les conteneurs principaux (cartes, modales) pour un look moderne et doux, et \textbf{4px} pour les éléments d'action (boutons, inputs) pour un aspect plus direct et structuré.
    \end{itemize}
\end{tcolorbox}

\textbf{Points de rupture (Breakpoints) :}
\begin{itemize}[label=\textcolor{GrillRed}{\faDesktop}]
    \item \textbf{Desktop (> 1024px) :} Affichage en grille multiple (ex: 3 cartes produits par ligne), menu de navigation complet visible.
    \item \textbf{Tablette (768px - 1024px) :} Optimisé pour le profil \textit{Restaurateur}. Affichage Kanban optimal.
    \item \textbf{Mobile (< 768px) :} Optimisé pour les profils \textit{Client} et \textit{Livreur}. Affichage en colonne unique, menu condensé sous format "Burger" (\faBars).
\end{itemize}

\section{10. Badges et Retours Visuels (Feedback)}
Les états des commandes et les alertes utilisent des codes couleurs universels pour une compréhension immédiate.

\begin{center}
\renewcommand{\arraystretch}{1.6}
\begin{tabular}{|m{4.5cm}|m{4cm}|m{6cm}|}
    \hline
    \rowcolor{CoalBlack} \textcolor{white}{\textbf{État / Statut}} & \textcolor{white}{\textbf{Style (Couleurs)}} & \textcolor{white}{\textbf{Utilisation typique}} \\
    \hline
    \textbf{En attente / Info} & Fond \textcolor{FryGold}{\textbf{Or}} (\#FFC107) & Nouvelle commande, en préparation. \\
    \hline
    \textbf{Succès / Validé} & Fond \textcolor{SuccessGreen}{\textbf{Vert}} (\#2E7D32) & Commande prête, paiement validé, livraison effectuée. \\
    \hline
    \textbf{Erreur / Annulée} & Fond \textcolor{GrillRed}{\textbf{Rouge}} (\#D32F2F) & Paiement refusé, rupture de stock, compte bloqué. \\
    \hline
    \textbf{Neutre / Historique} & Fond \textcolor{StoneGray}{\textbf{Gris}} (\#BDBDBD) & Commande passée, éléments désactivés. \\
    \hline
\end{tabular}
\end{center}

\section{11. Animations et Micro-interactions}
\begin{itemize}[label=\textcolor{CoalBlack}{\tiny\faCircle}, itemsep=0pt]
    \item \textbf{Survol (Hover) :} Au survol des boutons et des cartes produits, une transition douce (\texttt{transition: all 0.2s ease}) est appliquée avec un léger effet de zoom (\texttt{transform: scale(1.02)}) et l'accentuation de l'ombre portée pour inciter au clic.
    \item \textbf{Menus et Modales :} L'apparition des menus déroulants ou des fenêtres modales se fait en fondu. Pour les modales, l'arrière-plan de la page s'assombrit (overlay noir à 50\% d'opacité) pour concentrer l'attention de l'utilisateur.
\end{itemize}

\end{document}
